\documentclass[11pt]{article}
% Round-1 letter. Kept for the record: its PDF is frozen; recompiling it would pick up the current (round-2) numbers.
\usepackage[T1]{fontenc}
\usepackage[a4paper,margin=2.2cm]{geometry}
\usepackage{amsmath,amssymb}
\usepackage[per-mode=symbol]{siunitx}
\usepackage{enumitem}
\usepackage{xcolor}
\usepackage[hidelinks]{hyperref}
% Auto-generated by paper/make_paper_assets.py - do not edit by hand.

\newcommand{\YieldL}{1608}
\newcommand{\PRL}{0.819}
\newcommand{\CUFL}{18.35}
\newcommand{\POAL}{1963}
\newcommand{\TdayL}{39.5}
\newcommand{\TmaxL}{57.3}
\newcommand{\TriseL}{10.1}
\newcommand{\RHL}{74.0}
\newcommand{\DewL}{1{,}163}
\newcommand{\SRL}{0.965}
\newcommand{\DegBaseL}{0.30}
\newcommand{\DegTHL}{0.34}
\newcommand{\DegPIDL}{0.15}
\newcommand{\DegCorrL}{0.09}
\newcommand{\DegTotL}{0.88}
\newcommand{\TOWL}{3{,}055}
\newcommand{\RHencL}{63.5}
\newcommand{\LifeL}{36{,}233}
\newcommand{\CapL}{80.9}
\newcommand{\LCOEL}{2.83}
\newcommand{\LossIncidenceAngleL}{3.44}
\newcommand{\LossSpectralL}{-0.23}
\newcommand{\LossSoilingL}{3.67}
\newcommand{\LossDewCondensationL}{0.02}
\newcommand{\LossModuleTemperatureL}{6.12}
\newcommand{\LossDcWiringMismatchMotionLidL}{3.46}
\newcommand{\LossInverterAndClippingL}{2.03}
\newcommand{\LossAvailabilityL}{1.00}
\newcommand{\YieldF}{1634}
\newcommand{\PRF}{0.837}
\newcommand{\CUFF}{18.66}
\newcommand{\POAF}{1953}
\newcommand{\TdayF}{36.5}
\newcommand{\TmaxF}{50.0}
\newcommand{\TriseF}{7.3}
\newcommand{\RHF}{77.5}
\newcommand{\DewF}{1{,}092}
\newcommand{\SRF}{0.988}
\newcommand{\DegBaseF}{0.30}
\newcommand{\DegTHF}{0.34}
\newcommand{\DegPIDF}{0.15}
\newcommand{\DegCorrF}{0.05}
\newcommand{\DegTotF}{0.84}
\newcommand{\TOWF}{3{,}719}
\newcommand{\RHencF}{68.4}
\newcommand{\LifeF}{37{,}004}
\newcommand{\CapF}{81.7}
\newcommand{\LCOEF}{3.54}
\newcommand{\LossIncidenceAngleF}{3.47}
\newcommand{\LossSpectralF}{-0.24}
\newcommand{\LossSoilingF}{1.27}
\newcommand{\LossDewCondensationF}{0.04}
\newcommand{\LossModuleTemperatureF}{4.78}
\newcommand{\LossDcWiringMismatchMotionLidF}{4.62}
\newcommand{\LossInverterAndClippingF}{2.03}
\newcommand{\LossAvailabilityF}{1.50}
\newcommand{\YieldS}{1590}
\newcommand{\PRS}{0.814}
\newcommand{\CUFS}{18.15}
\newcommand{\POAS}{1953}
\newcommand{\TdayS}{36.2}
\newcommand{\TmaxS}{49.4}
\newcommand{\TriseS}{7.0}
\newcommand{\RHS}{79.7}
\newcommand{\DewS}{1{,}389}
\newcommand{\SRS}{0.969}
\newcommand{\DegBaseS}{0.30}
\newcommand{\DegTHS}{0.36}
\newcommand{\DegPIDS}{0.21}
\newcommand{\DegCorrS}{0.30}
\newcommand{\DegTotS}{1.16}
\newcommand{\TOWS}{4{,}288}
\newcommand{\RHencS}{70.4}
\newcommand{\LifeS}{34{,}677}
\newcommand{\CapS}{75.5}
\newcommand{\LCOES}{4.43}
\newcommand{\LossIncidenceAngleS}{3.47}
\newcommand{\LossSpectralS}{-0.25}
\newcommand{\LossSoilingS}{3.29}
\newcommand{\LossDewCondensationS}{0.05}
\newcommand{\LossModuleTemperatureS}{4.64}
\newcommand{\LossDcWiringMismatchMotionLidS}{4.91}
\newcommand{\LossInverterAndClippingS}{2.03}
\newcommand{\LossAvailabilityS}{2.00}
\newcommand{\YieldLM}{1608}
\newcommand{\PRLM}{0.819}
\newcommand{\CUFLM}{18.35}
\newcommand{\POALM}{1963}
\newcommand{\TdayLM}{39.5}
\newcommand{\TmaxLM}{57.3}
\newcommand{\TriseLM}{10.1}
\newcommand{\RHLM}{74.0}
\newcommand{\DewLM}{1{,}163}
\newcommand{\SRLM}{0.965}
\newcommand{\DegBaseLM}{0.30}
\newcommand{\DegTHLM}{0.34}
\newcommand{\DegPIDLM}{0.02}
\newcommand{\DegCorrLM}{0.04}
\newcommand{\DegTotLM}{0.70}
\newcommand{\TOWLM}{3{,}055}
\newcommand{\RHencLM}{63.5}
\newcommand{\LifeLM}{37{,}002}
\newcommand{\CapLM}{84.5}
\newcommand{\LCOELM}{2.90}
\newcommand{\LossIncidenceAngleLM}{3.44}
\newcommand{\LossSpectralLM}{-0.23}
\newcommand{\LossSoilingLM}{3.67}
\newcommand{\LossDewCondensationLM}{0.02}
\newcommand{\LossModuleTemperatureLM}{6.12}
\newcommand{\LossDcWiringMismatchMotionLidLM}{3.46}
\newcommand{\LossInverterAndClippingLM}{2.03}
\newcommand{\LossAvailabilityLM}{1.00}
\newcommand{\YieldFM}{1634}
\newcommand{\PRFM}{0.837}
\newcommand{\CUFFM}{18.66}
\newcommand{\POAFM}{1953}
\newcommand{\TdayFM}{36.5}
\newcommand{\TmaxFM}{50.0}
\newcommand{\TriseFM}{7.3}
\newcommand{\RHFM}{77.5}
\newcommand{\DewFM}{1{,}092}
\newcommand{\SRFM}{0.988}
\newcommand{\DegBaseFM}{0.30}
\newcommand{\DegTHFM}{0.34}
\newcommand{\DegPIDFM}{0.01}
\newcommand{\DegCorrFM}{0.03}
\newcommand{\DegTotFM}{0.68}
\newcommand{\TOWFM}{3{,}719}
\newcommand{\RHencFM}{68.4}
\newcommand{\LifeFM}{37{,}698}
\newcommand{\CapFM}{84.9}
\newcommand{\LCOEFM}{3.60}
\newcommand{\LossIncidenceAngleFM}{3.47}
\newcommand{\LossSpectralFM}{-0.24}
\newcommand{\LossSoilingFM}{1.27}
\newcommand{\LossDewCondensationFM}{0.04}
\newcommand{\LossModuleTemperatureFM}{4.78}
\newcommand{\LossDcWiringMismatchMotionLidFM}{4.62}
\newcommand{\LossInverterAndClippingFM}{2.03}
\newcommand{\LossAvailabilityFM}{1.50}
\newcommand{\YieldSM}{1590}
\newcommand{\PRSM}{0.814}
\newcommand{\CUFSM}{18.15}
\newcommand{\POASM}{1953}
\newcommand{\TdaySM}{36.2}
\newcommand{\TmaxSM}{49.4}
\newcommand{\TriseSM}{7.0}
\newcommand{\RHSM}{79.7}
\newcommand{\DewSM}{1{,}389}
\newcommand{\SRSM}{0.969}
\newcommand{\DegBaseSM}{0.30}
\newcommand{\DegTHSM}{0.36}
\newcommand{\DegPIDSM}{0.02}
\newcommand{\DegCorrSM}{0.15}
\newcommand{\DegTotSM}{0.83}
\newcommand{\TOWSM}{4{,}288}
\newcommand{\RHencSM}{70.4}
\newcommand{\LifeSM}{36{,}045}
\newcommand{\CapSM}{81.9}
\newcommand{\LCOESM}{4.43}
\newcommand{\LossIncidenceAngleSM}{3.47}
\newcommand{\LossSpectralSM}{-0.25}
\newcommand{\LossSoilingSM}{3.29}
\newcommand{\LossDewCondensationSM}{0.05}
\newcommand{\LossModuleTemperatureSM}{4.64}
\newcommand{\LossDcWiringMismatchMotionLidSM}{4.91}
\newcommand{\LossInverterAndClippingSM}{2.03}
\newcommand{\LossAvailabilitySM}{2.00}
\newcommand{\GainYieldF}{1.6}
\newcommand{\GainYieldS}{-1.1}
\newcommand{\GainLifeF}{2.1}
\newcommand{\GainLifeS}{-4.3}
\newcommand{\GainLifeSM}{3.9}
\newcommand{\dTdayF}{3.0}
\newcommand{\dTmaxF}{7.3}
\newcommand{\WaterSaved}{11{,}700}
\newcommand{\GHI}{1{,}889}
\newcommand{\Tair}{28.1}
\newcommand{\RHair}{74.0}
\newcommand{\Rain}{1{,}043}
\newcommand{\SensLifeMin}{32{,}498}
\newcommand{\SensLifeMax}{37{,}005}
\newcommand{\SensDegMin}{0.78}
\newcommand{\SensDegMax}{1.49}
\newcommand{\SensSaltSpan}{2{,}854}
\newcommand{\SensMoistSpan}{1{,}674}
\newcommand{\SensDegPerK}{0.167}
\newcommand{\SensTOWZero}{2{,}253}
\newcommand{\SensTOWMax}{5{,}547}
\newcommand{\SensYieldLoPct}{4.6}
\newcommand{\BreakEvenSalinity}{$<$0.25}
\newcommand{\NSeeds}{20}
\newcommand{\EnsYieldF}{1.71}
\newcommand{\EnsYieldSdF}{0.08}
\newcommand{\EnsLifeF}{2.17}
\newcommand{\EnsLifeSdF}{0.10}
\newcommand{\EnsDegF}{-0.039}
\newcommand{\EnsDegSdF}{0.003}
\newcommand{\EnsYieldS}{-1.08}
\newcommand{\EnsYieldSdS}{0.05}
\newcommand{\EnsLifeS}{-4.32}
\newcommand{\EnsLifeSdS}{0.06}
\newcommand{\EnsDegS}{+0.288}
\newcommand{\EnsDegSdS}{0.002}
\newcommand{\EnsYieldLandMin}{1{,}570}
\newcommand{\EnsYieldLandMax}{1{,}649}
\newcommand{\EnsFreshWins}{20}

% Auto-generated by paper/make_paper_assets.py (review additions) - do not edit by hand.

\newcommand{\NMC}{1{,}000}
\newcommand{\NParams}{28}
\newcommand{\MCYieldMedF}{1.2}
\newcommand{\MCYieldLoF}{$-$1.7}
\newcommand{\MCYieldHiF}{5.3}
\newcommand{\MCYieldPosF}{78}
\newcommand{\MCLifeMedF}{1.5}
\newcommand{\MCLifeLoF}{$-$3.6}
\newcommand{\MCLifeHiF}{6.3}
\newcommand{\MCLifePosF}{75}
\newcommand{\MCYieldMedS}{$-$1.5}
\newcommand{\MCYieldLoS}{$-$5.5}
\newcommand{\MCYieldHiS}{3.1}
\newcommand{\MCYieldPosS}{27}
\newcommand{\MCLifeMedS}{$-$5.1}
\newcommand{\MCLifeLoS}{$-$11.5}
\newcommand{\MCLifeHiS}{0.7}
\newcommand{\MCLifePosS}{4}
\newcommand{\MCYieldMedSM}{$-$1.5}
\newcommand{\MCYieldLoSM}{$-$5.5}
\newcommand{\MCYieldHiSM}{3.1}
\newcommand{\MCYieldPosSM}{27}
\newcommand{\MCLifeMedSM}{$-$1.6}
\newcommand{\MCLifeLoSM}{$-$7.0}
\newcommand{\MCLifeHiSM}{3.8}
\newcommand{\MCLifePosSM}{27}
\newcommand{\RCNameA}{Land dust deposition}
\newcommand{\RCValA}{0.45}
\newcommand{\RCNameB}{FPV $U_1$}
\newcommand{\RCValB}{0.43}
\newcommand{\RCNameC}{Saline salt deposition}
\newcommand{\RCValC}{$-$0.41}
\newcommand{\RCNameD}{FPV $U_0$}
\newcommand{\RCValD}{0.30}
\newcommand{\RCNameE}{Saline chloride deposition}
\newcommand{\RCValE}{$-$0.28}
\newcommand{\RCNameF}{Land $U_1$}
\newcommand{\RCValF}{$-$0.26}
\newcommand{\RCFNameA}{Land dust deposition}
\newcommand{\RCFValA}{0.59}
\newcommand{\RCFNameB}{FPV $U_1$}
\newcommand{\RCFValB}{0.49}
\newcommand{\RCFNameC}{FPV $U_0$}
\newcommand{\RCFValC}{0.35}
\newcommand{\PLcoeF}{99.5}
\newcommand{\PLcoeS}{100.0}
\newcommand{\OATBase}{$-$4.30}
\newcommand{\ClStarText}{none in 0--400}
\newcommand{\SaltStarYield}{35}
\newcommand{\DTdStar}{2.7}
\newcommand{\LifeSalineClZeroPct}{$-$0.9}
\newcommand{\LifeSalineSaltZeroPct}{$-$0.9}
\newcommand{\YieldSalineSaltZeroPct}{1.6}
\newcommand{\DegSalineClMax}{1.31}
\newcommand{\DegSalineClZero}{0.87}
\newcommand{\ClMax}{400}
\newcommand{\TOWFreshAtStar}{5{,}745}
\newcommand{\DewFreshMin}{1.64}
\newcommand{\DewFreshMax}{1.66}
\newcommand{\DewSalineMin}{$-$1.10}
\newcommand{\DewSalineMax}{$-$1.06}
\newcommand{\MitGGvsLand}{1.5}
\newcommand{\MitAllvsLand}{3.3}
\newcommand{\MitPIDRec}{2.2}
\newcommand{\MitSaltRec}{1.7}
\newcommand{\MitBothRec}{3.9}
\newcommand{\MitBothvsLand}{$-$0.5}
\newcommand{\CapexStarF}{34.3}
\newcommand{\CapexStarS}{29.1}
\newcommand{\LandPremF}{10.5}
\newcommand{\LandPremS}{23.6}
\newcommand{\WaterStar}{93}
\newcommand{\LCOEWaterZero}{3.54}
\newcommand{\LCOEWaterTen}{3.47}
\newcommand{\LCOEWaterTwenty}{3.39}
\newcommand{\LCOEWaterForty}{3.24}
\newcommand{\LCOEWaterEighty}{2.93}
\newcommand{\FinGapFMin}{0.58}
\newcommand{\FinGapFMax}{0.86}
\newcommand{\FinGapSMin}{1.36}
\newcommand{\FinGapSMax}{1.87}

\setlength{\parindent}{0pt}
\setlength{\parskip}{5pt}
\newcounter{cmt}
\newenvironment{comment}[1]{\refstepcounter{cmt}\par\medskip\noindent\textbf{Comment \thecmt{} (#1).}\itshape\ }{\par}
\newenvironment{response}{\par\noindent\textbf{Response.}\ }{\par}
\newcommand{\change}[1]{\par\noindent\textcolor{blue!50!black}{\textbf{Changes:} #1}\par}

\begin{document}
\begin{center}
{\Large\bfseries Response to Reviewers}\\[4pt]
Manuscript: ``Performance Analysis of Floating Solar Photovoltaic Systems for Humid and Saline Conditions: A Coupled
Thermal, Soiling and Degradation Assessment''\\
Revised title: ``Lifetime Performance and Degradation of Floating Photovoltaic Systems Under Humid and Saline
Conditions''
\end{center}

I thank the reviewer for a careful and constructive assessment. The revision addresses every comment. The largest
changes are:
\begin{enumerate}[nosep]
\item The IEC~61701/PID terminology is corrected throughout, and the mitigation factors are presented as explicit
assumptions (new Table~IV) and sampled in the uncertainty analysis.
\item A \NMC-sample Monte Carlo over \NParams{} uncertain parameters is added (new Section~III-I, Section~IV-H, Table~X,
Fig.~12). It is reported separately from the \NSeeds-year weather ensemble.
\item New threshold analyses (salt, chloride and dew-point rise; Fig.~9), a mitigation comparison (Table~VIII), economic
break-even values (Section~IV-I), a dew-loss sensitivity (Section~IV-C), a literature comparison table (Table~I) and a
weather-generator check (Table~III) are added.
\item Results are reframed as model-based scenario predictions throughout, and the scope is restricted explicitly to
module-level electrical degradation.
\end{enumerate}
One important outcome of the new uncertainty analysis is that the freshwater FPV lifetime advantage is \emph{not}
robust (positive in \MCLifePosF\% of samples), whereas the saline lifetime deficit is (negative in all but
\MCLifePosS\% of samples). The abstract, results and conclusion now state this. All numbers in the manuscript and in this
letter are generated by one script from the open-source model, so they are mutually consistent.

A correction to the generator also changed the baseline numbers slightly. The weather check requested by the reviewer
revealed that the synthetic year over-estimated mean wind speed by about \SI{0.5}{\metre\per\second} and mean RH by about
2.6 points relative to its input normals. Both biases have been removed (wind variability is now mean-preserving, and
the daily dew point is adjusted per month). The baseline saline lifetime deficit therefore changes from $-4.9$\% to
$\GainLifeS$\%, and the saline degradation rate from 1.24 to \DegTotS\%/yr. None of the qualitative conclusions changes.

\section*{Section 3: IEC 61701 associated with PID resistance}

\begin{comment}{terminology}
IEC~61701 is the salt-mist corrosion standard, not a PID qualification; IEC~TS~62804-1 is the PID standard. Correct the
terminology throughout.
\end{comment}
\begin{response}
Agreed; this was an error. The mitigated case is now described as ``PID-resistant module technology (qualified with
IEC~TS~62804-1:2025) combined with salt-mist hardening (qualified with IEC~61701:2020)''. The two effects are kept
separate in the model, the text, the figures and the code. The reference to IEC~TS~62804-1 is updated to the 2025
edition.
\end{response}
\change{Abstract; Sections~I, III-G, IV-E, V-B; Table~IV; captions of Fig.~7 and Tables~VII--VIII; bibliography; code
(\texttt{with\_mitigation} now has separate \texttt{pid} and \texttt{salt\_mist} switches).}

\begin{comment}{$\mu_{\mathrm{PID}}$ and $\mu_{\mathrm{C}}$}
The reductions $\mu_{\mathrm{PID}}=0.9$ and $\mu_{\mathrm{C}}=0.5$ are not justified by certification. Add a table of
standards and assumptions, or perform a sensitivity analysis.
\end{comment}
\begin{response}
Agreed. New Table~IV lists each mitigation, its qualification standard, how it enters the model and the assumed value
with its Monte Carlo range ($\mu_{\mathrm{PID}}$: 0.5--0.95; $\mu_{\mathrm{C}}$: 0.2--0.7). The text now states that
qualification standards define test procedures and do not establish a numerical reduction of field degradation. With
the factors sampled, the 95\% interval of the mitigated saline lifetime difference to land is [\MCLifeLoSM,
\MCLifeHiSM]\%. No quantitative literature value supports a specific reduction, so the factors are treated purely as
assumptions.
\end{response}
\change{Section~III-G and Table~IV (new); Section~IV-H and Table~X.}

\section*{Section 4: Environmental model largely assumed}

\begin{comment}{wording}
Present the 1.24\%/yr saline degradation as a scenario-model result under assumed environmental parameters.
\end{comment}
\begin{response}
Agreed. The results section now opens by stating that all results are model predictions under the baseline assumptions.
Absolute statements are rephrased (``Under the baseline assumptions, total degradation is \ldots''; ``the predicted
degradation rate''). Table~II is retitled ``Baseline Scenario Parameters (Assumptions)'', and the Monte Carlo interval
for the saline degradation rate is reported alongside the baseline value (Table~X).
\end{response}
\change{Abstract; Sections~II, IV (opening sentence), IV-D, VI.}

\section*{Section 5: Synthetic weather is not validation}

\begin{comment}{weather validation}
Validate against measured data or, at minimum, compare the synthetic weather statistics with measured data (GHI,
temperature, RH, rainfall, wind; RMSE, bias, correlation).
\end{comment}
\begin{response}
Measured hourly data for the site could not be obtained for this revision. New Table~III reports the MBE, RMSE and
correlation of \NSeeds{} generated years against the monthly normals the generator is built from. This check uncovered
and led to the correction of the wind and RH biases described above. The manuscript states explicitly that this is a
check of the generator against its inputs, not an external validation; that the RH match is exact by construction; and
that GHI (through the calibrated clearness index) has not been checked against measurements. The model accepts measured
hourly weather (\texttt{--weather}), and validation against measured meteorological and FPV data is named as the next
step. I acknowledge that this comment is only partly addressed.
\end{response}
\change{Section~III-A and Table~III (new); Section~V-C (Limitations); generator correction in
\texttt{fpv\_analysis/climate.py}.}

\section*{Section 6: Degradation model too empirical; module vs.\ system degradation}

\begin{comment}{scope}
Separate electrical/module degradation from FPV system degradation (structural fatigue, moorings, storm damage).
\end{comment}
\begin{response}
Agreed. A ``Scope'' paragraph at the end of Section~I states that the degradation results describe the electrical
degradation of the PV modules only, and that floating-system degradation is outside the model. Recent work on
structural corrosion and mechanical degradation of FPV components is cited. The limitations list the missing physics
(PID saturation and recovery, mechanism interactions, interconnect failure).
\end{response}
\change{Section~I (Scope); Section~V-C; ``module'' added to degradation labels in Tables~VII--VIII and Fig.~7.}

\section*{Section 7: Dew model}

\begin{comment}{4\% optical loss}
Justify the 4\% dew optical loss or show a sensitivity for 0, 2, 4 and 6\%.
\end{comment}
\begin{response}
The value is now labelled as an engineering assumption. New Section~IV-C shows that varying it over 0--6\% changes the
year-1 yield difference to land by less than 0.05 percentage point (freshwater: \DewFreshMin--\DewFreshMax\%; saline:
\DewSalineMin--\DewSalineMax\%), because condensate evaporates within the first hours after sunrise. The dew loss is
also sampled (0--6\%) in the Monte Carlo. Dew matters through time of wetness, not as an optical loss.
\end{response}
\change{Section~III-C; Section~IV-C (new).}

\section*{Section 8: Salt/soiling model}

\begin{comment}{salt deposition sensitivity}
Show salt deposition $\to$ first-year yield and salt deposition $\to$ lifetime energy.
\end{comment}
\begin{response}
Added as Fig.~9(a). Saline FPV out-yields land PV in year 1 only below about \SaltStarYield~mg/m$^2$/day of salt
deposition, and its 25-year energy stays below land at every salt level ($\LifeSalineSaltZeroPct$\% at zero salt). The
salt deposition, soiling coefficients and cementation factor are also sampled in the Monte Carlo; the saline salt
deposition is the third most influential parameter for the saline lifetime difference. Field evidence that sea-salt
deposits dominate on coastal FPV modules~[Ahmad \emph{et al.}, 2026] is now cited.
\end{response}
\change{Section~IV-F and Fig.~9 (new); Section~III-D.}

\section*{Section 10: Table III caption}

\begin{comment}{non-additive losses}
State that the sequential loss percentages are not directly additive.
\end{comment}
\begin{response}
Done. The table caption now reads ``Percentages Are Sequential Losses Relative to the Preceding Stage and Are Not
Directly Additive''.
\end{response}
\change{Table~VI caption; Fig.~4 caption.}

\section*{Section 11: Freshwater gain and parameter uncertainty}

\begin{comment}{weather vs.\ parameter uncertainty}
The $\pm0.08$\% from 20 weather realizations is stochastic weather uncertainty only; report parameter uncertainty for
the absolute yield.
\end{comment}
\begin{response}
Agreed. Weather variability and parameter uncertainty are now treated and reported separately (Section~IV-H). The
parameter uncertainty is much larger. The year-1 yield difference of freshwater FPV has a 95\% interval of
[\MCYieldLoF, \MCYieldHiF]\% (positive in \MCYieldPosF\% of samples), and its lifetime difference [\MCLifeLoF,
\MCLifeHiF]\% (\MCLifePosF\%). Table~X also gives intervals for the absolute yields, degradation rates and LCOEs. The
headline claim is softened accordingly: the freshwater advantage is a baseline result that depends mainly on the FPV
heat-loss coefficients and on how dusty the land alternative is.
\end{response}
\change{Abstract; Sections~III-I, IV-H (new); Table~X; Fig.~12; Sections~V and VI.}

\section*{Section 12: LCOE break-even}

\begin{comment}{break-even analysis}
Add break-even values for CAPEX, land cost, evaporation value and salt deposition.
\end{comment}
\begin{response}
Added (Section~IV-I). Because LCOE is linear in CAPEX and in the water value, the break-even values follow in closed
form:
\begin{itemize}[nosep]
\item Parity with land PV requires a CAPEX of \CapexStarF~INR/Wp for freshwater FPV and \CapexStarS~INR/Wp for saline
FPV.
\item Land PV would need a further \LandPremF{} or \LandPremS~INR/Wp of land-related cost to match them.
\item Freshwater parity requires a water value of about \WaterStar~INR/m$^3$.
\item No salt level makes saline FPV competitive at the assumed CAPEX.
\end{itemize}
In the Monte Carlo, FPV is more expensive than land PV in \PLcoeF\% (freshwater) and \PLcoeS\% (saline) of samples.
\end{response}
\change{Section~III-H (equation~(15) with water benefit); Section~IV-I (new).}

\section*{Section 13: Evaporation benefit}

\begin{comment}{monetize water}
Monetize the avoided evaporation and report the LCOE including the water benefit.
\end{comment}
\begin{response}
Done. The LCOE definition now includes an optional water benefit $B_y=c_wV_w$. At 10, 20 and 40~INR/m$^3$, the
freshwater LCOE falls from \LCOEWaterZero{} to \LCOEWaterTen, \LCOEWaterTwenty{} and \LCOEWaterForty~INR/kWh. The water
values are illustrative and are given as a range rather than as a single tariff.
\end{response}
\change{Equation~(15); Section~IV-I.}

\section*{Section 14: Literature positioning}

\begin{comment}{comparison table}
Add a systematic comparison with recent work in a feature table.
\end{comment}
\begin{response}
Added as Table~I, with 19 new references. The closest prior work is now discussed explicitly: Mannino \emph{et al.}
(2023) applied existing degradation models to offshore meteorological data, and Ajmal \emph{et al.} (2025) added
salinity to calibrated climatic degradation models for offshore modules. The novelty claim is narrowed to the coupling,
in one microclimate, of the yield benefits with the humidity-, salt- and PID-driven penalties, and to the
freshwater-versus-saline comparison with lifetime energy and cost. Recent FPV thermal measurements (Lindholm 2021;
D{\"o}renk{\"a}mper 2024; Nysted 2024; Wu 2024; F{\i}rat 2026) are now used to set the Monte Carlo ranges of the FPV
heat-loss coefficients, which include designs with no cooling advantage. PID studies on humidity and salt-induced surface
conduction (Hoffmann 2014; Hacke 2015; Luo 2017; Islam 2018; Bora 2021; Li 2025) now support the leakage term.
\end{response}
\change{Section~I and Table~I (new); Sections~III-C, III-F, V-A; bibliography.}

\section*{Section 16: Conclusions stronger than the evidence}

\begin{comment}{soften conclusions}
``Salinity is a stronger driver\ldots'' and ``module selection is the most effective mitigation'' should be presented as
model-based findings.
\end{comment}
\begin{response}
Done. The conclusion now reads: ``Under the baseline parameterization, salinity is the dominant driver of the predicted
lifetime penalty,'' and ``Module selection for PID and salt-mist resistance is the most effective modelled mitigation,
although its size rests on assumed reduction factors.'' It closes by stating that the conclusions are model-based and
call for field calibration.
\end{response}
\change{Section~VI; Section~V-B.}

\section*{Section 17: Figure~7}

\begin{comment}{model-derived}
Make the model-derived nature of the degradation components more prominent in Fig.~7.
\end{comment}
\begin{response}
The axis now reads ``Modelled degradation rate'', and the caption states ``Modelled module degradation rate by mechanism
(baseline)'' with a pointer to the assumed mitigation factors.
\end{response}

\section*{Section 18: Reproducibility}

\begin{comment}{reproducibility package}
Provide exact package versions, random seeds, configuration files, input weather, machine-readable parameters, unit tests
and a single command that reproduces everything.
\end{comment}
\begin{response}
All are now provided:
\begin{itemize}[nosep]
\item \texttt{requirements-lock.txt} pins the exact package versions (Python 3.11.15, TeX Live 2023).
\item \texttt{paper/generated/parameters.json} records every scenario parameter, the Monte Carlo ranges and all seeds
(baseline weather 42, ensemble 1--20, Monte Carlo 2024).
\item The synthetic weather is regenerated deterministically from these seeds.
\item 22 unit and regression tests cover, among other things, the reproducibility of the Monte Carlo and the agreement
of the generator with its normals.
\item \texttt{./reproduce.sh} regenerates every number, table and figure, the manuscript and this letter.
\end{itemize}
\end{response}

\section*{Sections 19--20: Recommended additional simulations}

\begin{comment}{Experiments A--D}
(A) Monte Carlo parameter uncertainty; (B) chloride threshold; (C) humidity threshold $\Delta T_d^*$; (D) mitigation
comparison.
\end{comment}
\begin{response}
All four have been added:
\begin{itemize}[nosep]
\item \textbf{(A)} \NMC{} samples over \NParams{} parameters, with medians, 95\% intervals, probabilities of a positive
difference and Spearman rank correlations (Table~X, Fig.~12).
\item \textbf{(B)} Saline degradation rises from \DegSalineClZero\%/yr at zero chloride to \DegSalineClMax\%/yr at
\ClMax~mg/m$^2$/day. No chloride level in this range restores lifetime parity with land: even at zero chloride the
deficit is $\LifeSalineClZeroPct$\%, because salt soiling, salt-assisted PID and the extra DC and availability losses
remain.
\item \textbf{(C)} Freshwater FPV loses its lifetime advantage above $\Delta T_d^*\approx\DTdStar$~K (Fig.~9b).
\item \textbf{(D)} Table~VIII compares six options for the saline plant: standard, PID-resistant, salt-mist-hardened,
both, both plus glass--glass, and all plus 7-day cleaning. Each is reported with recovery, comparison with land and
LCOE. Only with a lower-permeability package added does the saline plant exceed the land plant, under the stated
assumptions.
\end{itemize}
\end{response}
\change{Sections~III-G, III-I, IV-E, IV-F, IV-H; Tables~IV, VIII, X; Figs.~9 and 12.}

\section*{Sections 21--22: Title and abstract}

\begin{comment}{title and positioning}
Shorten the title; position the abstract as a model-based assessment.
\end{comment}
\begin{response}
The reviewer's second suggestion is adopted: ``Lifetime Performance and Degradation of Floating Photovoltaic Systems
Under Humid and Saline Conditions''. The abstract now opens with ``This study develops a coupled hourly simulation
framework to quantify\ldots'', reports the Monte Carlo intervals, and ends by stating that the results are model-based
scenario predictions that require field validation.
\end{response}

\section*{Remaining limitation}
The principal remaining weakness is the absence of measured validation data, both meteorological and FPV
module-temperature, performance and degradation data. This revision quantifies how much the conclusions depend on the
assumed parameters, but it cannot replace field data. The framework is ready to ingest measured hourly series, and
calibration against monitoring data from operating FPV plants in coastal Andhra Pradesh is planned.

\end{document}
